% !TeX spellcheck = en_US

\addscenariosection{1}{Tryb Starcia}{Miasto na Jeziorze}{\images/title.png}

\begin{multicols*}{2}

\textbf{Ten scenariusz można rozegrać tylko z rozszerzeniami: Twierdza, Zatoka, Wrota Żywiołów, Stretch Goals 2, Bitwy Morskie.}

\textbf{Autor:} Amnion

\textit{Na samotnej wyspie po środku wielkiego jeziora, magiczny eksperyment wyrwał się z pod kontroli. Wielka eksplozja doprowadziła do wysadzenia w powietrze całego miasta, niegdyś siedziby potężnej gildii magów. Pobliskie wody zapełniły się od gruzu, ciał mieszkańców i ich dobytku. Lokalny możnowładca, wynajął najemników zdolnych odnaleźć wszelkie wartościowe pozostałości zagubione wokół Miasta na Jeziorze. Za każdy odnaleziony skarb czeka was sowita nagroda. Może łączy was wspólny cel, ale tylko jednemu przypadną względy suwerena.}

\subsection*{\MakeUppercase{DŁUGOŚĆ SCENARIUSZA}}

Ten scenariusz rozgrywany jest w 16 rundach.

\subsection*{\MakeUppercase{Ustawienia Graczy}}

\textbf{Liczba Graczy:} 3

\textbf{Zasoby Początkowe:} 10 \svg{gold}

\textbf{Początkowy Przychód Graczy:} 10 \svg{gold}, 4 \svg{building_materials}, 1 \svg{valuables}

\textbf{Jednostki Początkowe:}
\begin{itemize}
  \item Dwie Garstki \svg{bronze} o najniższym koszcie rekrutowania. Następnie każdy gracz dobiera dwie losowe \svg{silver} Jednostki Neutralne, wybiera jedną i odrzuca drugą. Odrzucone w ten sposób jednostki neutralne wtasujcie do Talii Jednostek Neutralnych.
\end{itemize}

\textbf{Budynki Początkowe:} Siedlisko~\svg{bronze}, Budynek~\svg[12]{building_special_tent}


\columnbreak

\subsection*{\MakeUppercase{PRZYGOTOWANIE MAPY}}

Wymienione poniżej kafelki Mapy rozłóżcie zgodnie z przedstawionym dalej układem:
\begin{itemize}
  \item 3 × kafelek Początkowy (I)
  \item 6 × kafelek Daleki (II--III), każdy gracz powinien mieć dostęp do jednej Osady
  \item 6 × kafelek Bliski Morza (IV--V), trzy z nich muszą zawierać pole Puszka Pandory
  \item 6 × kafelek Bliski Podziemi (IV--V)
  \item 1 × kafelek Centralny Morza (VI--VII) z polem Świątynia Morza.
  \item 1 × kafelek Centralny Podziemi (VI--VII) z polem Losowe Miasto.
  \item 6 × żetonów Banków Stworzeń (II--III) 
  \item 3 × żetonów Banków Stworzeń (IV--V)

Dwa Kafelki Bliskie Morza posiadające pole z Puszką Pandory nie mogą ze sobą sąsiadować.
  
\end{itemize}

\subsection*{\MakeUppercase{WARUNEK ZWYCIĘSTWA}}
Na końcu ostatniej rundy gry wygrywa gracz który zdobył najwięcej kart Puszek Pandory. Jeżeli dwóch lub więcej graczy zebrała największą liczbę Puszek, gra kończy się remisem. Jeżeli na końcu którejkolwiek rundy jeden z graczy posiada osiem lub więcej Kart Puszki Pandory, zagrajcie jeszcze jedną turę. Na koniec tej tury wygrywa gracz który zdobył najwięcej kart Puszek Pandory.

\subsection*{\MakeUppercase{WARUNEK PRZEGRANEJ}}
Jeżeli na końcu ostatniej rundy całkowita licza kart Puszka Pandory wsród Graczy jest mniejsza niż 12, wszyscy przegrywają.
\newpage
\subsection*{\MakeUppercase{WYDARZENIA CZASOWE}}

\textbf{2. Runda:}
\begin{itemize}
  \item Każdy gracz może Przeszukać (2) \svg{artifact}.
\end{itemize}

\textbf{3. Runda:}
\begin{itemize}
  \item Rozpatrz Kartę Wydarzeń.
\end{itemize}

\textbf{4. Runda:}
\begin{itemize}
  \item Każdy z graczy dobiera Kartę Zwoju.
\end{itemize}

\textbf{6. Runda:}
\begin{itemize}
  \item Rozpatrz Kartę Wydarzeń.
\end{itemize}

\textbf{7. Runda:}
\begin{itemize}
  \item Każdy gracz rzuca i rozpatruje \svg{resource}.
\end{itemize}

\textbf{8. Runda:}
\begin{itemize}
  \item Każdy dostaje żeton Ulepszonej Umiejętności. Gracz z najmniejszą liczbą Kart Puszki Pandory dodatkowo dostaje \svg{movement}.
\end{itemize}

\textbf{9. Runda:}
\begin{itemize}
  \item Rozpatrz Kartę Wydarzeń.
\end{itemize}

\textbf{10. Runda:}
\begin{itemize}
  \item Zaczynając od tej rundy każda walka na polu Puszka Pandory ma zwiększony poziom trudności (np. Z Normalnego na Trudny). Efekt ten dotyczy także walki w Banku Stworzeń.
\end{itemize}

\textbf{12. Runda:}
\begin{itemize}
  \item Rozpatrz Kartę Wydarzeń.
\end{itemize}

\textbf{13. Runda:}
\begin{itemize}
  \item Gracz z najmniejszą liczbą Kart Puszki Pandory rzuca dwie \svg{resource} i rozpatruje jedną. Dodatkowo dostaje on \svg{morale_positive}.
\end{itemize}

\textbf{15. Runda:}
\begin{itemize}
  \item Rozpatrz Kartę Wydarzeń.
\end{itemize}


\subsection*{\MakeUppercase{DODATKOWE ZASADY}}

\begin{itemize}
    \item Walka z Neutralną Armią VII poziomu oraz walka na polu Puszka Pandory poziomu nie może zostać pominięta.
    \item Po dobraniu Karty Puszki Pandory i rozpatrzeniu jej efektu odłóż ją na bok tworząc stos Kart Puszek Pandory Gracza.
    \item Co Rundę Astrologów możesz wybrać jedną posiadaną Kartę Puszki Pandory i rozpatrzyć ją ponownie. Karta może być rozpatrzona w ten sposób tylko raz na grę. Efekty zwiększające Poziom Doświadczenia Gracza nie mogą być rozpatrzone ponownie.
    \item Po wygranej walce na Polu Banku Stworzeń (II-III) zamiast rozpatrywać normalną nagrodę, dobierz Kartę Puszki Pandory. W przypadku Banków (VI-V) nagrodą jest zarówno bazowy efekt jak i Karta Puszki Pandory.
    \item Nagrodą za pokonanie Głównego Bohatera przeciwnika jest jedna z jego Kart Puszki Pandory. Jeżeli dany gracz nie posiada takiej karty, zmuszony jest oddać 8 \svg{gold} lub równowartość tej kwoty w zasobach. 
    \item Podczas bitwy każdy z Graczy może przekazać Kartę Puszki Pandory aby dokonać odwrotu do miasta zachowując swoją armię.
    \item Zajęcie wrogiego miasta umożliwia zabranie wybranej Karty Puszki Pandory gracza do którego należało miasto. W wypadku w którym przejęcie miasta było rezultatem wygranego oblężenia, atakujący może przejąć tylko jedną Kartę Puszki Pandory. Miasto nie może zostać zajęte i może być zaatakowane przez każdego z wrogich graczy raz na grę.
    \item Podczas rozpatrywania pola Sanktuarium, Gracz może Przeszukać(2) swoją talię.
    \item Podczas rozpatrywania pola z Obeliskiem, Gracz może dobrać po jednej karcie z Talii Jednostek Neutralnych odpowiadających posiadanym przez niego Siedliskom, a następnie zrekrutować jedną z nich.
    \item Podczas rozpatrywania pola Świątynia Morza, zamień jedną \svg{azure} Jednostkę Neutralną na Hydry oraz dołóż Jednostkę Neutralną Haspidy. Zamiast standardowej nagrody otrzymujesz 20\svg{gold}, Przeszukaj(4)~\svg{artifact} oraz zamień jedną ze swoich Kart Statystyk na Kartę Ulepszonej Statystyki.
    \item Podczas rozpatrywania pola Losowe Miasto. Zamiast standardowej nagrody możesz \svg{pay} aby Zrekrutować i Ulepszyć wybrane jednostki z frakcji Losowego Miasta za połowę ceny. Dodatkowo jedną ze swoich Umiejętności na Kartę Ulepszonej Umiejętności.
    \item Podczas rozpatrywania pola z Neutralną Armią VII, dodatkowo dobierz Kartę Puszki Pandory.
\end{itemize}

\end{multicols*}



\noindent
\includegraphics[
  width=0.8\paperwidth,
  height=0.5\textheight,
  keepaspectratio
]{\maps/owcy_potwor_w_3p.png}

\begin{center}
  \captionof{figure}{\textbf{3-PLAYER SCENARIO}}
\end{center}
